\documentclass[10pt,oneside]{report}
\usepackage{soc2style}
\begin{document}
\doctitle{DevSecOps and Application Security Mapping}{GitHub security controls, SAST, DAST, IAST and SDLC assurance}{SOC2-06}
\section*{Operating convention and boundaries}
An AICPA criterion identifies an outcome against which controls may be evaluated. The selected point-of-focus headings identify relevant characteristics; they are not discrete certification checklist items. Every owner assignment, technological mechanism and audit test in this suite is a proposed example to be tailored to the system description, commitments, risk profile, period, and service organization responsibilities.

\begin{itemize}
\item Always establish system boundaries, the in-scope trust services categories, service commitments and system requirements first.
\item For a security-only engagement, use the common criteria. For additional trust services categories, include the corresponding additional criteria, subject to properly justified applicability.
\item Document complementary user entity controls and subservice organizations; do not treat vendor assurance reports as replacements for management ownership.
\item Privacy points of focus marked [P] or [C] describe typical processor or controller relevance; assess actual responsibilities rather than treating the annotations as absolute.
\end{itemize}

\chapter{Assurance boundaries for application security}
The AICPA source does not mandate GitHub, CodeQL, GitHub Advanced Security (GHAS), Dependabot, DAST, IAST, SARIF, SBOMs or any particular DevSecOps product. These are contemporary implementation options to help management meet relevant criteria. Tool-based claims must be verified against deployed configuration, license entitlements, repositories, branches and time-bound control operation.

\subsection{Controls versus scanner features}
\begin{itemize}
\item Capability: a vendor or scanner can identify a class of potential findings; it may be unavailable, misconfigured or outside scope.
\item Control activity: an assigned, measurable process evaluates and acts on those outputs, with exceptions, monitoring and traceable results.
\item Assurance evidence: source records demonstrate that the control operated for the relevant full population, period and objective.
\item Residual risk: even a deployed detector has false negatives, alert noise, blind spots, unmodeled architectures and unmonitored repositories.
\end{itemize}

\chapter{AppSec control-to-criteria mapping}
\section{Source control organization governance}
Suggested criteria alignment: CC1.3 CC2.1 CC6.1 CC6.3. The final mapping depends on the risks and controls in the scoped system.

\begin{engineeringbox}
Repository inventory, ownership, administrator review, ruleset exemptions
\end{engineeringbox}
Candidate evidence: Organization security exports, membership, privileged role reviews.

Validation: derive an authoritative population of affected applications, repositories or services; identify required check executions; confirm missing or bypassed results are represented; trace findings to triage and verified remediation or an approved, time-bounded exception.

\section{Branch protections and review}
Suggested criteria alignment: CC5.3 CC6.3 CC8.1. The final mapping depends on the risks and controls in the scoped system.

\begin{engineeringbox}
Mandatory PR review, approval independence, up-to-date checks, emergency change traceability
\end{engineeringbox}
Candidate evidence: Rule exports, branch bypass audit, PR approval and merge records.

Validation: derive an authoritative population of affected applications, repositories or services; identify required check executions; confirm missing or bypassed results are represented; trace findings to triage and verified remediation or an approved, time-bounded exception.

\section{CodeQL and SAST}
Suggested criteria alignment: CC7.1 CC7.3 CC8.1. The final mapping depends on the risks and controls in the scoped system.

\begin{engineeringbox}
In-scope language/config analysis, alert handling, remediation checks before affected merges
\end{engineeringbox}
Candidate evidence: SARIF, alert metadata, analysis configuration, successful/failed check runs.

Validation: derive an authoritative population of affected applications, repositories or services; identify required check executions; confirm missing or bypassed results are represented; trace findings to triage and verified remediation or an approved, time-bounded exception.

\section{Secrets prevention, discovery and response}
Suggested criteria alignment: CC6.1 CC6.6 CC7.1 CC7.3. The final mapping depends on the risks and controls in the scoped system.

\begin{engineeringbox}
Secrets scanning, protected token storage, validation, rotation or revocation
\end{engineeringbox}
Candidate evidence: Alerts, identity provider activity, rotation receipts, exception decision records.

Validation: derive an authoritative population of affected applications, repositories or services; identify required check executions; confirm missing or bypassed results are represented; trace findings to triage and verified remediation or an approved, time-bounded exception.

\section{Dependency and supply-chain hygiene}
Suggested criteria alignment: CC3.2 CC7.1 CC9.2 CC8.1. The final mapping depends on the risks and controls in the scoped system.

\begin{engineeringbox}
Dependency inventory, advisories, reachability review, approved build artifact and version pinning
\end{engineeringbox}
Candidate evidence: Dependency manifests, SBOM, build provenance, patch and accept tickets.

Validation: derive an authoritative population of affected applications, repositories or services; identify required check executions; confirm missing or bypassed results are represented; trace findings to triage and verified remediation or an approved, time-bounded exception.

\section{DAST and dynamic verification}
Suggested criteria alignment: CC7.1 CC7.3 CC8.1. The final mapping depends on the risks and controls in the scoped system.

\begin{engineeringbox}
Risk-based scans of deployed exposed services with approved scope and safe test plans
\end{engineeringbox}
Candidate evidence: Target inventory, test runs, verified findings and remediation retests.

Validation: derive an authoritative population of affected applications, repositories or services; identify required check executions; confirm missing or bypassed results are represented; trace findings to triage and verified remediation or an approved, time-bounded exception.

\section{IAST and runtime observability}
Suggested criteria alignment: CC7.1 CC7.2 CC7.3. The final mapping depends on the risks and controls in the scoped system.

\begin{engineeringbox}
Runtime evidence for executed vulnerable paths, context, deployment linkage and error trends
\end{engineeringbox}
Candidate evidence: Instrumentation coverage, request traces, validation tickets.

Validation: derive an authoritative population of affected applications, repositories or services; identify required check executions; confirm missing or bypassed results are represented; trace findings to triage and verified remediation or an approved, time-bounded exception.

\section{Deployment security and secrets in CI}
Suggested criteria alignment: CC6.1 CC6.3 CC6.7 CC8.1. The final mapping depends on the risks and controls in the scoped system.

\begin{engineeringbox}
Short-lived credentials, least privilege, environment protection and auditable artifact release
\end{engineeringbox}
Candidate evidence: Workflow identity policy, pipeline logs, deployment approvals, signature/provenance.

Validation: derive an authoritative population of affected applications, repositories or services; identify required check executions; confirm missing or bypassed results are represented; trace findings to triage and verified remediation or an approved, time-bounded exception.

\section{Cloud and container misconfiguration}
Suggested criteria alignment: CC2.1 CC6.6 CC7.1 CC7.2. The final mapping depends on the risks and controls in the scoped system.

\begin{engineeringbox}
Inventory reconciliation, config drift tracking, hardened workloads and admission rules
\end{engineeringbox}
Candidate evidence: IaC scan outputs, cloud drift events, scanner coverage, change history.

Validation: derive an authoritative population of affected applications, repositories or services; identify required check executions; confirm missing or bypassed results are represented; trace findings to triage and verified remediation or an approved, time-bounded exception.

\section{Vulnerability exception governance}
Suggested criteria alignment: CC3.2 CC4.2 CC7.3. The final mapping depends on the risks and controls in the scoped system.

\begin{engineeringbox}
Rational disposition taxonomy, business justification, compensating controls, expiry and approval
\end{engineeringbox}
Candidate evidence: Issue decision history, risk acceptance record and retest results.

Validation: derive an authoritative population of affected applications, repositories or services; identify required check executions; confirm missing or bypassed results are represented; trace findings to triage and verified remediation or an approved, time-bounded exception.

\section{Production incident and response}
Suggested criteria alignment: CC2.2 CC2.3 CC7.3 CC7.4 CC7.5. The final mapping depends on the risks and controls in the scoped system.

\begin{engineeringbox}
Secure alert routing, incident command, customer communication decisions and validated restoration
\end{engineeringbox}
Candidate evidence: Security incident evidence packet, tabletop record, PR/release remediation.

Validation: derive an authoritative population of affected applications, repositories or services; identify required check executions; confirm missing or bypassed results are represented; trace findings to triage and verified remediation or an approved, time-bounded exception.

\section{Data protection in SDLC}
Suggested criteria alignment: CC6.7 C1.1 C1.2 P3.1 P4.1 P4.2 P4.3. The final mapping depends on the risks and controls in the scoped system.

\begin{engineeringbox}
Data minimization, approved libraries, input/output validation, privacy-by-design review and lifecycle management
\end{engineeringbox}
Candidate evidence: Threat model, data-flow diagram, test fixtures, retention/deletion records.

Validation: derive an authoritative population of affected applications, repositories or services; identify required check executions; confirm missing or bypassed results are represented; trace findings to triage and verified remediation or an approved, time-bounded exception.

\chapter{CodeQL assessment operating procedure}
\begin{enumerate}
\item Define asset and repository population from authoritative organization inventory, including active and archived repositories, forks where relevant, and external build dependencies.
\item Determine CodeQL eligibility and analysis language coverage for each repository; distinguish no code, unsupported language, disabled configuration, failed run and successful run with no findings.
\item Review repository triggers for default branches, protected branches and pull requests that represent the system change process; identify branches where no required checks apply.
\item Extract alert records with tool, rule, severity, alert state, created/fixed/dismissed timestamps, locations, alert URL and last scan outcome.
\item Deduplicate alerts by appropriate stable identities while retaining per-repository and per-commit provenance; do not aggregate distinct occurrences into one defect silently.
\item Triage exploitability by code context, source-to-sink path, sanitizers, effective access boundaries and deployed configuration. Never claim CodeQL determines exploitability of every finding.
\item For remediation, identify the fix commit, confirm a fresh applicable scan and link the fixed result to the authorized merge/deployment pathway.
\item Produce decision records for non-fix outcomes with documented justification, approved risk owner if risk is accepted, expiration and periodic review.
\item Monitor coverage and backlog trends together; do not reward alert count reduction caused by disabling scans or broad dismissals.
\end{enumerate}

\subsection{Illustrative CodeQL dashboard fields}
\begin{lstlisting}[language={},caption={Suggested analytics schema, not an API response contract}]
organization, repository, default_branch, archived, language
codeql_enabled, supported_language, last_analysis_at
last_analysis_commit, last_analysis_result
pull_request_required_check, bypass_enabled
alert_id, rule_id, severity, alert_state
alert_created_at, alert_fixed_at, alert_dismissed_at
triage_outcome, approver_role, remediation_ticket
exception_expiration, evidence_export_timestamp
\end{lstlisting}
\chapter{Secret scanning triage decision record}
\begin{enumerate}
\item Classify the detected material: service credential, key material, access token, benign public identifier, test fixture or unknown.
\item Establish source, accessible exposure, live credential validity, scopes, ability to authenticate and risk from historical copies.
\item When valid or sensitive, rotate or revoke promptly, investigate misuse and document effective remediation time.
\item When non-sensitive or non-exploitable, record concrete validation facts, technical owner, source code context and bounded rationale; do not use an unsubstantiated blanket false-positive designation.
\item When intentionally retained for a test environment, validate that it is not production-capable, is scoped, monitored and approved; repeat assessment if conditions change.
\end{enumerate}

\subsection{Sample disposition comment patterns}
\begin{itemize}
\item Remediated: "Credential revoked and replaced; old token failed authentication at verification time; associated incident/risk ticket linked."
\item Validated non-secret: "String matches a public client identifier with no authentication privilege; tested usage and configuration; owner validated classification."
\item Exception: "Remediation deferred for an approved reason, compensating controls in place, accountable owner, review date and expiry recorded."
\end{itemize}

Disposition names and governance meanings here are examples. Use the actual organization policy and available platform disposition types; do not attribute a preferred suppression taxonomy to the AICPA.

\chapter{Enterprise metrics without perverse incentives}
\begin{itemize}
\item Repository coverage = repositories with an applicable, successful analysis / repositories eligible and in scope; report exclusions separately.
\item Effective gate coverage = protected-branch merge paths with enforced, verified checks / applicable merge paths; include exceptions and bypasses.
\item Finding age = measurement time minus initial finding time, grouped by risk, affected exposure and remediation ownership.
\item Time to validated remediation = verification time minus triage or detection time; preserve disposition history.
\item Secret response quality = cases with validity check, rotation/revocation decision and verification / eligible secret cases.
\item Vulnerability risk posture = confirmed, potentially exploitable, unknown, and accepted risk classes; never equate the count of SAST findings with the inventory of vulnerabilities.
\end{itemize}

\chapter{AppSec audit packet acceptance criteria}
\begin{enumerate}
\item The scope identifies all intended organizations, repos, applications, components and environments.
\item Collection date, analysis rules, scan triggers and relevant branches are documented.
\item Operational sample ties finding or change to the exact commit, PR, reviewer and check run.
\item Missing/incomplete scanning is tracked separately from no findings.
\item Risk acceptances are attributed and time-bound; fixed outcomes are backed by retesting.
\item Evidence is reproducible through stable IDs, immutable exports and controlled retrieval.
\end{enumerate}


\chapter*{Source and editorial protocol}
\addcontentsline{toc}{chapter}{Source and editorial protocol}
This independent document is grounded in the user-supplied \emph{AICPA TSP Section 100: 2017 Trust Services Criteria for Security, Availability, Processing Integrity, Confidentiality, and Privacy (With Revised Points of Focus---2022)}, 75-page PDF, especially paragraphs .01--.29, the criterion table on pp.~14--67, and the glossary. Criterion wording is quoted when labeled as such. Selected point-of-focus labels are aids to interpretation and may be incomplete. Engineering recommendations, role assignments, tools, tests, dates, scorecards, and example frequencies are author-developed examples and are not AICPA requirements. The 2022 revisions changed points of focus, not the 2017 criteria (source p.~3). The source states that an assessment of each point of focus is not required (p.~5). This document does not constitute a CPA attestation or legal advice.
\end{document}
